\chapter{Цветосохраняющий составной циклический родитель}
\label{ch:v7-color-preserving-composite-cycle-parent-gate}

\section{Ориентированное слово цикла}

Полный gauge-подъём запретил ненулевые фундаментальные значения двух
цветных рёбер. Однако укоренённый шестикромочный цикл является замкнутым
gauge-инвариантным словом
\begin{equation}
 \mathcal W_C
 =D_{Q_Lu_R}D_{u_RX_L}D_{X_Le_R}
  D_{e_RL_L}D_{L_LY_R}D_{Y_RQ_L}.
 \label{eq:v7-composite-cycle-full-word}
\end{equation}
Цветовые и абелевы действия на соседних концах сокращаются внутри
композиции, поэтому $\Tr\mathcal W_C$ является gauge-синглетом. Это
стандартная причина, по которой спектральные действия на колчанах
раскладываются по замкнутым путям \cite{v7-composite-cycle-quiver}.

После фиксации двух старых корневых рёбер $H_{15}$ переменная часть имеет
вид
\begin{equation}
 C(z)=
 z_{u_RX_L}z_{X_Le_R}z_{L_LY_R}z_{Y_RQ_L}.
 \label{eq:v7-composite-cycle-new-word}
\end{equation}
Два множителя
\begin{equation}
 z_{u_RX_L}\in\overline{\mathbf3}_{-5/3},
 \qquad
 z_{Y_RQ_L}\in\mathbf3_{2/3}
 \label{eq:v7-composite-cycle-colored-factors}
\end{equation}
цветно сопряжены; оставшиеся электрослабые множители замыкают полный
гиперзаряд слова.

\section{Классический логический барьер}

Для обычной конфигурации полей произведение ненулево только если ненулев
каждый множитель:
\begin{equation}
 C(z)\ne0
 \quad\Longrightarrow\quad
 z_{u_RX_L}\ne0,\qquad z_{Y_RQ_L}\ne0.
 \label{eq:v7-composite-cycle-nonzero-factors}
\end{equation}
Следовательно, ненулевой классический циклический параметр порядка снова
требует два фундаментальных triplet VEV и не сохраняет $SU(3)_c$.
Переименование $C$ в составной синглет не меняет это тождество.

Можно ввести независимое поле $\Sigma$ и ограничение
\begin{equation}
 \Sigma=C(z).
 \label{eq:v7-composite-cycle-auxiliary-constraint}
\end{equation}
Но тогда $\Sigma\ne0$ при точном классическом ограничении также влечёт
$C(z)\ne0$. Если же ограничение удалить, $\Sigma$ является новым
постулированным синглетом, а не производным циклом.

\section{Цикл не запускает нуль}

После заморозки старых рёбер $C(z)$ имеет степень четыре по новым полям.
Поэтому
\begin{equation}
 \nabla C(0)=0,\qquad
 \Hess_0 C=0.
 \label{eq:v7-composite-cycle-zero-hessian}
\end{equation}
Любой аналитический потенциал, первая зависимость которого от цикла имеет
вид $\Re C$, $|C|^2$ или более высокой степени, не создаёт квадратичной
отрицательной моды в нуле.

В частности, при положительных цветосохраняющих массах
\begin{equation}
 V(z)=m^2\sum_{j=1}^{4}|z_j|^2
 -\kappa\bigl(C(z)+\overline{C(z)}\bigr)+O(|z|^4),
 \qquad m^2>0,
 \label{eq:v7-composite-cycle-massive-potential}
\end{equation}
получаем
\begin{equation}
 \Hess_0V=2m^2I>0,
 \label{eq:v7-composite-cycle-positive-origin}
\end{equation}
независимо от $\kappa$. Циклический член может участвовать в удалённом
переходе первого рода, но любой классический минимум с $C\ne0$ снова имеет
ненулевые цветные фундаментальные поля.

\section{Граница квантового композита}

В gauge-инвариантной квантовой формулировке физические параметры порядка
должны быть gauge-инвариантными; ненулевое ожидание неинвариантного
локального поля не является наблюдаемым критерием
\cite{v7-composite-cycle-elitzur}. Но возможность
\begin{equation}
 \langle z_{u_RX_L}\rangle
 =\langle z_{Y_RQ_L}\rangle=0,
 \qquad
 \langle C(z)\rangle\ne0
 \label{eq:v7-composite-cycle-quantum-condensate}
\end{equation}
является утверждением о квантовой мере и связанных корреляторах, а не о
минимуме текущего классического Hodge-функционала.

Поэтому честный следующий маршрут должен оставить цветные мосты массивными
и интегрировать их как виртуальные поля. Детерминант или дополнение Шура
может породить gauge-инвариантный эффективный цикл при нулевых одноточечных
средних; такой механизм ещё не вычислен.

\begin{theorem}[Запрет классического составного обхода]
Укоренённое шестикромочное слово является точным gauge-синглетом. Но после
фиксации старого фона оно имеет четвёртую степень по новым полям, нулевые
градиент и гессиан в начале и ненулево только при ненулевых двух цветных
множителях. Следовательно, оно не запускает цветосохраняющий классический
вакуум и не устраняет предыдущий цветовой запрет.
\end{theorem}

\begin{nogocriterion}[Запрет переименования произведения в конденсат]
Нельзя объявлять $C(z)$ независимым цветосинглетным полем, сохраняя
одновременно точное равенство $\Sigma=C(z)$ и нулевые фундаментальные
цветные поля. Допустимое переоткрытие требует явно выведенной квантовой
меры, детерминанта или дополнения Шура, которое создаёт составной коррелятор
без triplet VEV.
\end{nogocriterion}

\begin{protocol}
\begin{enumerate}
 \item полный цикл gauge-инвариантен: \textbf{да};
 \item степень по новым полям после укоренения: \textbf{$4$};
 \item градиент цикла в нуле: \textbf{$0$};
 \item гессиан цикла в нуле: \textbf{$0$};
 \item ненулевой классический цикл без цветных множителей: \textbf{нет};
 \item положительные массы сохраняют локально устойчивый нуль: \textbf{да};
 \item независимое $\Sigma$ выведено текущим действием: \textbf{нет};
 \item квантовый составной конденсат проверен: \textbf{нет};
 \item классический составной маршрут: \textbf{закрыт};
 \item следующий гейт: \textbf{виртуальный цветной мост и дополнение Шура}.
\end{enumerate}
\end{protocol}

Машинный сертификат сохранён в
\path{s2t_v7_color_preserving_composite_cycle_parent_gate_results.json}.

\begin{thebibliography}{9}
\bibitem{v7-composite-cycle-quiver}
C.~I.~P\'erez-S\'anchez,
\emph{The Spectral Action on Quivers}, arXiv:2401.03705.
\bibitem{v7-composite-cycle-elitzur}
S.~Elitzur,
\emph{Impossibility of spontaneously breaking local symmetries},
Physical Review D 12 (1975), 3978--3982.
\end{thebibliography}